\section{Roteiro lâmina a lâmina}
\label{sec:roteiro}

A Tabela~\ref{tab:roteiro} detalha, para cada lâmina, o recurso visual, os dados apresentados e a origem desses dados. Valores projetados são sempre identificados como hipóteses em validação.

\begin{table}[htb]
\centering
\caption{Recursos visuais, dados e fontes por lâmina}
\label{tab:roteiro}
\fontsizetable
\begin{tabularx}{\textwidth}{L{0.8cm} L{3.4cm} Y L{3.4cm}}
\toprule
\textbf{Slide} & \textbf{Recurso visual} & \textbf{Dados e conteúdo} & \textbf{Fonte}\\
\midrule
1 & Ícone e identidade visual; nomes da equipe & Nome, \textit{slogan}, integrantes, disciplina. & Orientação da Atividade XII\\
2 & Ilustração em ícones da persona e do pet; linha do tempo de três momentos & Situação, conflito e resolução; consulta de R\$~150 a R\$~350. & Entrevistas (Atividade VI); levantamento de preços (Atividade IX)\\
3 & Números em destaque (46,1\%; 19,3\%; 28\%) & Lares com cães e com gatos; lares em que nem todos os animais estavam vacinados contra a raiva. & IBGE, PNS 2019 \cite{ibge2019pns,correio2020}\\
4 & Quatro cartões com ícones & Prontuário e calendário de vacinas; triagem orientativa; rede solidária; agendamento. & Aplicativo \cite{cassino2026repo}\\
5 & Diagrama de dois lados & Tutores de classes C, D e E; clínicas e ONGs. & Atividade VI; \citeonline{rochet2003}\\
6 & Três colunas (tutor, clínica, parceiro) & Benefícios por lado da plataforma. & Atividade IX\\
7 & Gráfico em anel e quatro valores & Composição da receita estimada no mês 12 (R\$~4.421,50). & Atividade IX (estimativa do grupo)\\
8 & Matriz comparativa & Plano de saúde pet, aplicativo de registro e PetCuida. & Atividade IX\\
9 & Números em destaque e linha do tempo & Cinco sessões de teste; dois formatos testáveis; economia unitária. & Atividades VI e IX\\
10 & Linha do tempo em quatro marcos & Piloto, teste de preço, formalização, segurança e LGPD. & Atividades IX, X e XIII\\
11 & Quatro cartões de integrantes & Nome e papel de cada integrante; pedido. & Atividade XI\\
\bottomrule
\end{tabularx}
\fonteautores
\end{table}

\subsection{Dados centrais do pitch}

Segundo a Pesquisa Nacional de Saúde de 2019, do IBGE, 46,1\% dos domicílios brasileiros tinham ao menos um cachorro (33,8 milhões de lares) e 19,3\% ao menos um gato (14,1 milhões) \cite{correio2020}. Em 72,0\% dos domicílios com cães ou gatos, todos os animais haviam sido vacinados contra a raiva \cite{conexaoto2020}; o complemento indica que, em cerca de 28\% desses domicílios, nem todos os animais estavam vacinados. Esse dado sustenta a mensagem de que a prevenção ainda é lacuna relevante.

Quanto ao modelo, a estimativa do grupo para o mês 12 é de 1.500 usuários ativos mensais, 150 atendimentos intermediados, 75 assinantes do PetCuida+ e 10 clínicas parceiras, com receita total de R\$~4.421,50, dos quais R\$~1.500,00 dependem de convênios e patrocínio ainda não contratados (Atividade IX). Na apresentação, esses valores são rotulados como estimativa, e a dependência dos convênios é explicitada.
